La integración multisensorial, el proceso mediante el cual el cerebro combina información de distintas modalidades sensoriales, constituye un área fundamental de la neurociencia computacional. Los modelos existentes presentan una brecha significativa: aquellos que capturan la arquitectura funcional de la integración multisensorial carecen de dinámicas corticales realistas, mientras que los modelos con dinámicas biológicamente plausibles están limitados al procesamiento unimodal. Este trabajo aborda dicha brecha mediante la implementación del modelo de \citet{echeveste2020cortical} dentro del framework Scikit-NeuroMSI, una plataforma de código abierto en Python para la modelización neurocomputacional de la integración multisensorial. Dicho modelo demuestra que redes neuronales con restricciones biológicas pueden realizar inferencia probabilística mediante muestreo estocástico, exhibiendo propiedades emergentes como oscilaciones gamma y variabilidad modulada por el estímulo. La reimplementación, fue validada alcanzando equivalencia numérica exacta con errores del orden de $10^{-15}$. Como contribuciones adicionales, se desarrollaron módulos para la gestión de parámetros pre-entrenados y para modelos generativos. Los resultados posicionan a esta implementación como base técnica para futuras investigaciones que combinen dinámicas corticales realistas con integración multisensorial.